\documentclass[10pt,a4paper]{amsart}
\usepackage[margin=19mm]{geometry}
\usepackage{amsmath,amssymb,mathtools}
\usepackage[scr=rsfs,cal=euler]{mathalfa}
\usepackage{xfrac}
\usepackage[unicode,hidelinks,bookmarks=false]{hyperref}
\newcommand{\ie}{i.e.\ }
\newcommand{\cf}{cf.\ }
\newcommand{\resp}{respectively\ }
\newcommand{\Subsection}{\subsection}
\providecommand{\nobreakdash}{\nobreak\hbox{-}}
\setlength{\emergencystretch}{2em}
\multlinegap0pt
\usepackage{enumerate}
\usepackage{amssymb}
\usepackage{mathtools}
\usepackage{float}
\newtheorem{theorem}{Theorem}
\newtheorem{proposition}{Proposition}
\newtheorem{lemma}{Lemma}
\usepackage{cleveref}
\crefname{theorem}{Theorem}{Theorems}
\crefname{proposition}{Proposition}{Propositions}
\crefname{lemma}{Lemma}{Lemmas}
\newcommand{\CC}{\mathbb{C}}
\newcommand{\ZZ}{\mathbb{Z}}
\newcommand{\cA}{\mathcal{A}}
\newcommand{\gen}[1]{\left\langle #1\right\rangle}
\title{On the Frobenius Number of Quotients of Numerical Semigroups}
\author{Feihu Liu}
\date{Source: arXiv:2607.23076v1 (CC BY 4.0)}
\begin{document}
\maketitle

\noindent\emph{Source and licence.} On the Frobenius Number of Quotients of Numerical Semigroups. Version arXiv:2607.23076v1 (CC BY 4.0), distributed by arXiv under the Creative Commons Attribution 4.0 International licence, http://creativecommons.org/licenses/by/4.0/, as recorded in the arXiv licence metadata for this exact version and shown on the abstract page https://arxiv.org/abs/2607.23076v1. Full licence notice: \texttt{licenses/arxiv-2607.23076-CC-BY-4.0.txt}.\par\smallskip

\noindent\textit{Editorial scope.} Counted unit: the article's theorem thm:main (source0.tex lines 231--238). The article's complete original proof of it is delivered with it (source0.tex lines 734--799, one \texttt{proof} environment of 66 lines, which the article places away from the statement). No mathematics is written, repaired or re-derived: the statement and the proof are the article's own lines, in the article's own notation, and no retained line is substituted or re-flowed.  Retained uncounted as the source-local context the proof needs: lines 477--495; lines 602--614; lines 664--677; lines 713--715; lines 720--722.  Nothing was omitted from any retained span, so the package carries no omission record.  No retained line contains a citation, so the wrapper carries no bibliography and no bibliography entry.  No retained line contains a figure environment, an image-inclusion command or drawing code, so no image is delivered and the package declares no asset dependency.  Editorial formatting only, in addition: an AMS \texttt{amsart} preamble that restates the article's own declarations and macros used by the retained lines, namely the environments \texttt{theorem}, \texttt{proposition}, \texttt{lemma} on separate counters, together with the standard packages those lines need.  The retained environments print in this document as Theorem 1, Proposition 1, Lemma 1 in order of appearance, whereas the article prints the same items with the numbering of its own full document.  The complete native source of the article as distributed in the arXiv e-print for this exact version is bundled unchanged as \texttt{sources/205995/source0.tex}.
\begin{proposition}\label{prop:explicit-branch}
Let $p>2$ be an odd integer, and let $k$ be an integer satisfying
\begin{equation}\label{eq:k-range}
  2\le k\le\frac{p+1}{2}.
\end{equation}
Let $a,b\in\ZZ_{>0}$ satisfy $p<a<b$, $\gcd(a,b)=1$,
\begin{equation}\label{eq:branch-congruences}
a\equiv1\pmod p, \qquad b\equiv p-k+1\pmod p,
\end{equation}
and
\begin{equation}\label{eq:branch-slope}
  p-k<\frac{b}{a}<p-k+1.
\end{equation}
Then
\begin{equation}\label{eq:explicit-branch}
  g\!\left(\frac{\gen{a,b}}{p}\right)
  =\frac{ab-(k-1)a-2b}{p}.
\end{equation}
\end{proposition}

\begin{lemma}\label{lem:approximation}
Let $p>2$ be a prime, let $2\le k\le\frac{p+1}{2}$, and let $\alpha\in(p-k,p-k+1)$.
Then there exist sequences of positive integers $(a_n)_{n\ge1}$ and
$(b_n)_{n\ge1}$ such that, for every $n$,
\begin{enumerate}
  \item $a_n$ is prime and $a_n\equiv1\pmod p$;
  \item $b_n\equiv p-k+1\pmod p$;
  \item $p<a_n<b_n$, $\gcd(a_n,b_n)=1$, and $p\nmid b_n$;
  \item $p-k<\frac{b_n}{a_n}<p-k+1$;
  \item $\lim_{n\to\infty}\frac{b_n}{a_n}=\alpha$.
\end{enumerate}
In particular, $(a_n,b_n,p)\in\cA$ for every $n$.
\end{lemma}

\begin{lemma}\label{lem:asymptotic-vanishing}
Let $I\subset\mathbb R$ be a nonempty open interval, and let
$G\in\CC[X_1,X_2]$.  Assume that for every $\alpha\in I$ there exists a
sequence $(x_n,y_n)_{n\ge1}$ of real pairs such that
\begin{equation}\label{eq:asymptotic-hypotheses}
  x_n\longrightarrow+\infty,
  \quad
  \frac{y_n}{x_n}\longrightarrow\alpha,\quad \text{as}\quad n\longrightarrow\infty;
  \qquad
  G(x_n,y_n)=0
  \quad\text{for all }n.
\end{equation}
Then $G$ is the zero polynomial.
\end{lemma}

\begin{lemma}\label{lem:root-bound}
Let $K$ be a field and let $0\neq H\in K[Y]$ have degree $d$.  Then $H$ has at most $d$ distinct roots in $K$.
\end{lemma}

\begin{lemma}\label{lem:specialization}
Let $F\in\CC[X_1,X_2,X_3,Y]$ be nonzero.  Then $F(X_1,X_2,q,Y)\in\CC[X_1,X_2,Y]$ is nonzero for all but finitely many $q\in\CC$.
\end{lemma}

\begin{theorem}\label{thm:main}
There is no nonzero polynomial $F\in\CC[X_1,X_2,X_3,Y]$ such that
\begin{equation}\label{eq:main-relation}
  F\left(a,b,p,
  g\!\left(\frac{\gen{a,b}}{p}\right)\right)=0
\end{equation}
for every $(a,b,p)\in\cA$.
\end{theorem}

\begin{proof}[Proof of Theorem~\ref{thm:main}]
Assume, for contradiction, that a nonzero polynomial $F\in\CC[X_1,X_2,X_3,Y]$ satisfies~\cref{eq:main-relation} for every $(a,b,p)\in\cA$.  Let $d:=\deg_Y F$. By Lemma~\ref{lem:specialization}, the specialization
$F(X_1,X_2,q,Y)$ is nonzero for all but finitely many complex numbers $q$.
There are infinitely many primes, so we may choose an odd prime $q$ such that
$H_q(X_1,X_2,Y) :=F(X_1,X_2,q,Y)\neq 0$.
and $\frac{q-1}{2}>d$.

For each integer
\begin{equation}\label{eq:k-values-main}
  k\in\left\{2,3,\ldots,\frac{q+1}{2}\right\},
\end{equation}
define
\begin{align*}
\Phi_{q,k}(X_1,X_2):=\frac{X_1X_2-(k-1)X_1-2X_2}{q}\in\CC[X_1,X_2].
\end{align*}
We shall prove that
\begin{equation}\label{eq:H-vanishes-branch}
  H_q\bigl(X_1,X_2,\Phi_{q,k}(X_1,X_2)\bigr)\equiv0
\end{equation}
for every $k$ in~\cref{eq:k-values-main}.

Fix such a $k$, and set $G_{q,k}(X_1,X_2):=H_q\bigl(X_1,X_2,\Phi_{q,k}(X_1,X_2)\bigr)$.
Let $I_{q,k}:=(q-k,q-k+1)$.
Fix an arbitrary $\alpha\in I_{q,k}$.  By Lemma~\ref{lem:approximation}, there exist sequences $(a_n)$ and $(b_n)$ such
that
\[(a_n,b_n,q)\in\cA,\qquad\frac{b_n}{a_n}\longrightarrow\alpha\quad \text{as}\quad n\longrightarrow \infty,
\]
and all the hypotheses of Proposition~\ref{prop:explicit-branch} hold with
$p=q$.  Therefore
\begin{align*}
  g\!\left(\frac{\gen{a_n,b_n}}{q}\right) =\Phi_{q,k}(a_n,b_n)
\end{align*}
for every $n$.

Since $(a_n,b_n,q)\in\cA$, the assumed polynomial relation gives
\begin{align*}
0&=F\left(a_n,b_n,q,g\!\left(\frac{\gen{a_n,b_n}}{q}\right)\right)
=H_q\left(a_n,b_n,g\!\left(\frac{\gen{a_n,b_n}}{q}\right)\right)
\\& =H_q\bigl(a_n,b_n,\Phi_{q,k}(a_n,b_n)\bigr)
=G_{q,k}(a_n,b_n).
\end{align*}
Also $a_n\to\infty$ because the $a_n$ are an increasing sequence of primes,
and $b_n/a_n\to\alpha$.  Since this construction works for every
$\alpha\in I_{q,k}$, Lemma~\ref{lem:asymptotic-vanishing} implies $G_{q,k}\equiv0$.
This proves~\cref{eq:H-vanishes-branch}.

Now let $K:=\CC(X_1,X_2)$ be the field of rational functions in $X_1$ and $X_2$.  By
\cref{eq:H-vanishes-branch}, each $\Phi_{q,k}$ is a root of the nonzero
polynomial $H_q\in K[Y]$.  These roots are pairwise distinct.  Indeed, if
$k\neq\ell$, then
\begin{align*}
\Phi_{q,k}-\Phi_{q,\ell}=\frac{-(k-1)X_1+(\ell-1)X_1}{q}=\frac{\ell-k}{q}X_1,
\end{align*}
which is a nonzero element of $K$.

The set~\cref{eq:k-values-main} contains exactly $\frac{q-1}{2}$ integers.  Hence $H_q$, viewed as a nonzero polynomial in $K[Y]$, has at least $(q-1)/2$ distinct roots.  By Lemma~\ref{lem:root-bound},
\begin{equation}\label{eq:degree-lower}
  \deg_Y H_q\ge\frac{q-1}{2}.
\end{equation}
On the other hand, specialization in $X_3$ cannot increase the degree in
$Y$, so
\begin{equation}\label{eq:degree-upper}
  \deg_Y H_q\le\deg_Y F=d.
\end{equation}
Combining~\cref{eq:degree-lower} and~\cref{eq:degree-upper} gives $d\ge\frac{q-1}{2}$, contradicting~$\frac{q-1}{2}>d$.  This contradiction proves that no such nonzero polynomial $F$ exists.
\end{proof}

\end{document}
